\documentclass[11pt,a4paper]{article}
\usepackage[utf8]{inputenc}
\usepackage[T1]{fontenc}
\usepackage[portuguese]{babel}
\usepackage{xcolor}
\usepackage{hyperref}
\usepackage{geometry}
\usepackage{listings}
\usepackage{tcolorbox}
\geometry{margin=2.5cm}
\definecolor{primary}{RGB}{41,128,185}
\definecolor{secondary}{RGB}{44,62,80}
\definecolor{codebg}{RGB}{240,240,240}
\lstset{backgroundcolor=\color{codebg},basicstyle=\ttfamily\small,breaklines=true,frame=single}
\title{\color{secondary}\Huge\textbf{Solução: Exercício 2.1 - Services e comunicação interna}}
\author{\color{primary}DevOps com Kubernetes -- 2026}
\date{}
\begin{document}
\maketitle


\textbf{Guia de Solução:}

\subsection*{\color{secondary}Passo a Passo}

1. Crie um Deployment com label \texttt{app: web} e três réplicas.
2. Crie um Service ClusterIP selecionando \texttt{app: web}.
3. Use \texttt{kubectl run curl --rm -it --image=curlimages/curl -- sh} e acesse \texttt{http://web-svc}.
4. Valide com \texttt{kubectl get endpoints web-svc} se o Service encontrou os Pods.

\end{document}
